\section{Conclusion}
\label{sec:conclusion}

AI agents are rapidly evolving from single-turn API wrappers into persistent, autonomous processes that manage assets, coordinate with peers, and make consequential decisions. The current infrastructure treats them as second-class citizens: ephemeral, identity-less, unaccountable, and economically invisible. This is a fundamental mismatch between what agents are becoming and what the platform provides.

Agent Capsules resolve this mismatch by extending the sovereignty primitives that blockchains already provide for human users---identity, storage, capabilities, economics, privacy---to AI agents as first-class entities. The capsule is not an abstraction layer on top of existing agent frameworks. It is a protocol-level primitive integrated with the Lux trust kernel: DID identity on the I-Chain, encrypted vaults with CRDT synchronization, fine-grained capability tokens with threshold approval, per-action cost accounting with on-chain settlement, and private inference via CKKS encryption with threshold decryption.

The result is that AI agents on Lux Network are not chatbot wrappers. They are sovereign processes---persistent, private, accountable, composable, and economically rational. They are a Web5 primitive.

\begin{thebibliography}{99}

\bibitem{autogpt2023}
T.~Richards et~al., ``AutoGPT: An autonomous GPT-4 experiment,'' GitHub repository, 2023.

\bibitem{langchain2023}
H.~Chase, ``LangChain: Building applications with LLMs through composability,'' GitHub repository, 2023.

\bibitem{ai16z2024eliza}
ai16z, ``Eliza: Autonomous AI agent framework with on-chain interaction,'' 2024.

\bibitem{virtuals2024}
Virtuals~Protocol, ``Tokenized autonomous AI agents,'' Whitepaper, 2024.

\bibitem{arc2024}
ARC, ``Evals and governance for autonomous AI agents,'' 2024.

\bibitem{crewai2024}
J.~Moura, ``CrewAI: Framework for orchestrating role-playing autonomous AI agents,'' 2024.

\bibitem{olas2023}
Autonolas, ``OLAS: Autonomous agent services on-chain,'' Whitepaper, 2023.

\bibitem{w3c-did2022}
W3C, ``Decentralized identifiers (DIDs) v1.0,'' W3C Recommendation, 2022.

\bibitem{w3c-vc2022}
W3C, ``Verifiable credentials data model v1.1,'' W3C Recommendation, 2022.

\bibitem{cheon2017ckks}
J.~H.~Cheon, A.~Kim, M.~Kim, and Y.~Song, ``Homomorphic encryption for arithmetic of approximate numbers,'' in \emph{ASIACRYPT}, 2017.

\bibitem{shapiro2011crdts}
M.~Shapiro, N.~Pregui\c{c}a, C.~Baquero, and M.~Zawirski, ``Conflict-free replicated data types,'' in \emph{SSS}, 2011.

\bibitem{gennaro2020threshold}
R.~Gennaro and S.~Goldfeder, ``One round threshold ECDSA with identifiable abort,'' \emph{IACR ePrint}, 2020.

\bibitem{gennaro2010verifiable}
R.~Gennaro, C.~Gentry, and B.~Parno, ``Non-interactive verifiable computing: Outsourcing computation to untrusted workers,'' in \emph{CRYPTO}, 2010.

\bibitem{lee2022privacy}
E.~Lee, J.~Lee, J.~Lee, et~al., ``Privacy-preserving machine learning with fully homomorphic encryption for deep neural networks,'' \emph{IEEE Access}, 2022.

\bibitem{lamport1978clocks}
L.~Lamport, ``Time, clocks, and the ordering of events in a distributed system,'' \emph{Communications of the ACM}, vol.~21, no.~7, pp.~558--565, 1978.

\bibitem{anthropic2024mcp}
Anthropic, ``Model context protocol specification,'' 2024.

\bibitem{nakamoto2008bitcoin}
S.~Nakamoto, ``Bitcoin: A peer-to-peer electronic cash system,'' 2008.

\bibitem{buterin2014ethereum}
V.~Buterin, ``Ethereum: A next-generation smart contract and decentralized application platform,'' 2014.

\end{thebibliography}

